\documentclass[11pt]{article}
\usepackage{ghb}

% Restored from the original manuscript on 2026-10-30.
% Source: /home/nick/Downloads/Coding-20240916T060449Z-001/Coding/Nick/misc/Writing/cycle.tex
% Source SHA-256: b9a19250422d4121d1a1fc50d0d12c552432b6737f6ccbce93aeb8efc85965d8
% Date is the archived source modification date, corroborated by the
% accompanying PDF creation date; original composition/publication is unknown.
% Prose retained verbatim; paragraph spacing adapted to the blog.
% Decoded original Windows-1252 punctuation as UTF-8.
% Original editorial note retained here, outside the story:
% All needs careful review. E.g. wondered a LITTLE if .. LITTLE .. LITTLE -- OOPS!! fix this and other problems

\title{On the Seriousness of Eating Eggs}
\date{2009-02-15}
\summary{Short fiction, from a February 2009 manuscript.}

\begin{document}

Patrick rose early, exhausted himself, and returned home contented and a bit too hot at dinner time. Sometimes in the afternoon his mother would call him in from the road. Wearing a pretty summer dress, flushed with the summer romance of a happy marriage, she handed him cool drinks. Usually Patrick would approach his mother with mixed feelings. She was disturbing his summer joys, and yet he was young enough that her smell or the sound of her voice still comforted all his anxieties and calmed his little heart.

``Come and have something to drink pet''.

He drained the drink and handed it back with a boyish smile. Each time, he ran to the street again, relaxed and buoyed up by the touch of his mother.

He wondered a little if his mother's life was not a little dull, and with a little anxiety sometimes had feelings of pity rise inside him. He once asked a friend:

``What do you think it is, when you're a grown up?''\\ ``I think it's when have a feeling that, that you feel like someone you don't know.''

Years later, the summer glow of his mother was only a memory. Patrick's brain still chanted his friend's words to him. He thought they must have been meaningless when he heard them, and he couldn't remember now when it was they began to cycle around and around in his head.

He remembered his friends, how they had been asleep. Now he was woken, as he looked around he saw some people still having pleasant dreams. Some were having nightmares. Occasionally he saw another wide-eyed face, pale and worried. He half thought they might be awake too, but he was afraid to ask. Most of the time he pretended to be asleep too.

Patrick was over fifty now, he still lived where his parents had. For a period of nearly fifteen years after secondary school he had lived independently though. Some women had even considered him for a mate. In the end though, terrible, warped shafts of light shone through him where he had left spaces.

Until recent years he had kept the house tidy enough. He had gone to the office parties. He'd made an effort to hold onto normality, he had engaged with the world. He read a little, he watched TV. Just to be sociable, he had even looked at sports results. He sated colleagues by inquiring with exaggerated ineptitude about their sport of choice, if they had one.

Remembering his childhood, it sometimes seemed that it had just been smudged into today. He remembered studying at school and his first job. He remembered thinking then that one day life would begin. Now, the terrible dullness of his spirit terrified him. At times he felt the echoes of his childhood happiness. The memories seemed to bounce inside him violently. Panicked, they dented him from inside, as if making desperately to leave his unhappy frame.

At night, lying alert in bed, he scrunched his eyes and trembled in disbelief at the horror of his ruined morale.

Since he lost his job the life-long pressure of stayed misery and nonsense finally had a steady flow. He was profoundly depressed, but rationally so, he thought. He was disgusted by human interaction, he longed to embrace all of humanity with pure love, but knew that really, he did not exist.

The house was filthy now, Patrick rarely went outside, and rarely even ate. Patrick was filthy too, and he especially avoided washing his hands. It was a challenge to himself, to demonstrate his rebellion against evolution and good sense. Shameful evolution.

One day, he struggled a little against his disengagement from normality, and tried to tidy up a bit. There were piles of books and newspapers all over the upstairs landing. He was carrying a pile of heavy books downstairs to throw them out and he slipped. He didn't quite fall down the stairs but shifted heavily onto his left ankle. It gave a spiteful little grind.

He dropped the books.

It felt as if someone had piled a thousand bricks on top of his ankle and foot. He pulled up his trouser leg a little and saw his hideously thin ankle, there was lots of blood on his sock, a bone had broken the skin.

Patrick saw the wound as a challenge to himself. He decided he must ignore it. He knew pain was a trick to manipulate his behaviour. He would not be an automaton. He sat on the stairs for some hours as he felt his shoe fill with blood. Eventually, he dozed off.

Apparently, life and hope were not fully unwound inside him. He woke. As his eyes focused he saw the doctor's kind face, there was even a pretty nurse, and she smiled warmly.

``Goodness me, why wouldn't you take care of yourself? You were a sad state coming in here after that accident. Aren’t you lucky to have a sister like that to bring you in?''

Patrick nodded weakly, pretty women, pointless as they were, still made him feel sad.

``Well don't worry now, you're looking great and when you're out of here you'll be straight back on your feet. Even under your fingernails, I've never seen so much dirt. No but don't worry now. Now it's clean, now…''

Patrick freed his right hand from under the duvet. It looked nothing like his hand. The nails had been neatly trimmed and the fingers were matte and clean. He had seen fingers like these on other people.

``Why have you cleaned even my hands?''

He was still slightly sedated. His words slurred and he wasn't sure if the nurse understood him. She smiled with a slight tilt of her head. He relaxed his neck, his head sinking into the pillow now. He began to drift off and as his eyes darkened he repeated to himself,

``The darkness lives under his fingernails, he keeps it next to his skin.''

When Patrick was discharged, he found his house lying in wait. Its terrifying familiarity caused his mind to rush, sweeping itself into the absolute silence of all of life’s pointlessness. With resolve, his spirit did not break as he crossed its threshold. The house had prepared. With a flourish, it laid down its gauntlet. He looked and saw the place, at the foot of the stairs, where he had met his father up one night. This was the place he first saw his father’s steel give in to elderly confusion. That night his father had not quite found the bathroom. He had come downstairs and then impatiently turned to his returning son and asked him where it was. This was a year or so after his mother had passed on. Mercifully, his father wasn’t too slow to follow. That had been a relief.

Patrick moved through this cursed place and into the kitchen and that decade-old relief still disgusted him.

He had made it to the kitchen now, and as he turned and saw the dusty curtains, he suddenly grew tired. Without warning, and without effort, he tired of the grinding sadness. He looked around the dirt atop the relics of his life so far. He felt the terrible pressure in his mind gently ease.

He opened the fridge. Stocked by his sister. He saw the little chick, dead from the cold. He saw its path now, it had hatched! A supermarket egg had hatched. The chick had tried bravely to find a life, but it had hatched in a fridge. The tender little chick never realized it was born into a fridge.

Today Patrick would try again, and commit his lunatic sadness to the langoliers. He laughed! He had just thought of eating an egg. Now, toast with marmalade.

\end{document}
